\textbf{1. Kiến thức cần nhớ}
\begin{itemize}
\item Phương trình đường tròn $(C)$ tâm $I(a; b)$, bán kính $R$: $\boxed{(x-a)^2 + (y-b)^2 = R^2}$[cite: 1].
\item Phương trình khai triển: $x^2 + y^2 - 2ax - 2by + c = 0$[cite: 1] là đường tròn nếu $a^2 + b^2 - c > 0$[cite: 1]. Khi đó, tâm $I(a; b)$ và $R = \sqrt{a^2 + b^2 - c}$[cite: 1].
\item Phương trình tiếp tuyến của $(C)$ tại điểm $M_0(x_0; y_0)$ thuộc $(C)$: $(a - x_0)(x - x_0) + (b - y_0)(y - y_0) = 0$[cite: 1].
\item Điều kiện tiếp xúc: Đường thẳng $\Delta$ tiếp xúc với đường tròn $(C)$ tâm $I$, bán kính $R$ $\Leftrightarrow d(I, \Delta) = R$[cite: 1].
\end{itemize}

\textbf{2. Phương pháp giải các dạng}

%% DANG 0H9?4-1
\textbf{Dạng 1. Tìm tâm, bán kính và điều kiện là đường tròn}
\begin{enumerate}
\item Từ phương trình $(x-a)^2 + (y-b)^2 = R^2$, đọc ngay tâm $I(a; b)$ và bán kính $R$.
\item Từ phương trình $x^2 + y^2 - 2ax - 2by + c = 0$, chia hệ số của $x, y$ cho $-2$ để được tọa độ tâm $I(a; b)$.
\item Kiểm tra điều kiện $a^2 + b^2 - c > 0$. Nếu thỏa mãn thì $R = \sqrt{a^2 + b^2 - c}$[cite: 1].
\end{enumerate}
Lưu ý: Hệ số của $x^2$ và $y^2$ phải bằng nhau (thường là $1$). Nếu hệ số là $2, 3...$ thì phải chia cả hai vế cho hệ số đó trước khi tìm $a, b, c$.

%% DANG 0H9?4-2
\textbf{Dạng 2. Phương trình đường tròn}
\begin{enumerate}
\item Tìm tọa độ tâm $I(a; b)$ và bán kính $R$.
\item Thay vào dạng $(x-a)^2 + (y-b)^2 = R^2$[cite: 1].
\item Nếu bài cho đi qua $3$ điểm, gọi dạng $x^2 + y^2 - 2ax - 2by + c = 0$[cite: 1], thay tọa độ $3$ điểm vào lập hệ $3$ phương trình ẩn $a, b, c$ và bấm máy giải hệ.
\end{enumerate}
Lưu ý: Đường tròn đường kính $AB$ có tâm là trung điểm $AB$ và $R = \frac{AB}{2}$[cite: 1]. Đường tròn tâm $I$ tiếp xúc đường thẳng $\Delta$ thì $R = d(I, \Delta)$[cite: 1]. Đừng quên bình phương $R$ khi viết phương trình[cite: 1].

%% DANG 0H9?4-3
\textbf{Dạng 3. Phương trình tiếp tuyến của đường tròn}
\begin{enumerate}
\item \textbf{Tại một điểm $M_0$ thuộc đường tròn:} Tiếp tuyến đi qua $M_0$ và có VTPT $\vec{n} = \overrightarrow{IM_0}$[cite: 1]. Viết phương trình $a(x-x_0) + b(y-y_0) = 0$.
\item \textbf{Biết phương/hệ số góc:} Gọi phương trình tiếp tuyến có dạng tổng quát (chứa tham số).
\item Ép điều kiện $d(I, \Delta) = R$ để giải tìm tham số.
\end{enumerate}
Lưu ý: "Tiếp tuyến \textbf{tại} $M$" khác "tiếp tuyến \textbf{kẻ từ} $M$". Nếu tiếp tuyến kẻ từ một điểm ngoài đường tròn, phải dùng điều kiện $d(I, \Delta) = R$.

%% DANG 0H9?4-4
\textbf{Dạng 4. Vị trí tương đối liên quan đường tròn}
\begin{enumerate}
\item Tính khoảng cách từ tâm $I$ đến đường thẳng $\Delta$ (hoặc khoảng cách $2$ tâm $I_1, I_2$).
\item So sánh khoảng cách đó với bán kính $R$:
  - $d < R$: Cắt nhau (tại 2 điểm).
  - $d = R$: Tiếp xúc.
  - $d > R$: Không giao nhau.
\end{enumerate}
Lưu ý: Để tìm tọa độ giao điểm, giải hệ phương trình gồm phương trình đường thẳng và đường tròn (thường dùng phương pháp thế).

%% DANG 0H9?4-5
\textbf{Dạng 5. Toán tổng hợp đường thẳng và đường tròn}
\begin{enumerate}
\item Nếu điểm $M$ thuộc đường thẳng, tham số hóa tọa độ $M$ theo phương trình đường thẳng[cite: 1].
\item Nếu liên quan đến dây cung $AB$ của đường tròn tâm $I$, gọi $H$ là trung điểm $AB$, dùng Pitago: $R^2 = IH^2 + \left(\frac{AB}{2}\right)^2$.
\item Lập phương trình theo yêu cầu bài toán để giải tham số.
\end{enumerate}
Lưu ý: Luôn vẽ hình nháp để nhìn ra mối liên hệ vuông góc, Pitago hoặc khoảng cách.

%% DANG 0H9?4-6
\textbf{Dạng 6. Bài toán dùng cho tam giác, tứ giác}
\begin{enumerate}
\item Đường tròn ngoại tiếp tam giác là đường tròn đi qua $3$ đỉnh (Dạng 2). Tâm là giao $3$ đường trung trực.
\item Đường tròn nội tiếp tam giác tiếp xúc $3$ cạnh.
\item Kết hợp các yếu tố hình học phẳng (góc, khoảng cách) đưa vào hệ tọa độ.
\end{enumerate}
Lưu ý: Tọa độ tâm đường tròn ngoại tiếp có thể tìm nhanh bằng cách gọi $I(x; y)$ rồi giải hệ $IA = IB$ và $IA = IC$.

%% DANG 0H9?4-7
\textbf{Dạng 7. Bài toán thực tế, PP tọa độ hóa}
\begin{enumerate}
\item Gắn hệ trục $Oxy$ vào mô hình thực tế (ví dụ: tâm đường cong/vòm cầu làm gốc $O$).
\item Tìm tọa độ các điểm quan trọng (đỉnh, tâm, điểm đi qua) dựa vào kích thước thực tế[cite: 1].
\item Viết phương trình đường tròn và giải các yêu cầu (tìm chiều cao, khoảng cách...).
\end{enumerate}
Lưu ý: Cần chọn đơn vị trên hệ trục tọa độ khớp với đơn vị đo lường trong thực tế (ví dụ: $1$ đơn vị $= 1\text{m}$).
